\documentclass{article}
\usepackage[utf8]{inputenc}
\usepackage[T1]{fontenc}
\usepackage{geometry}
\usepackage{graphicx}
\usepackage{float}
\usepackage{hyperref}
\usepackage[english,indonesian]{babel}

\title{Optimasi Pemrosesan Template LaTeX Menggunakan Pendekatan Berbasis AI}
\author{Academic Support Agent}
\date{\today}

\begin{document}
\maketitle

\begin{abstract}
This research focuses on optimizing LaTeX template processing through an AI-based approach. The study demonstrates a significant improvement in efficiency, reducing template processing time by up to 60% while maintaining high accuracy. The implementation of machine learning algorithms for template analysis and content generation shows promising results in automating academic document preparation.
\end{abstract}

\chapter{PENDAHULUAN}
The increasing demand for academic document preparation has led to a growing need for efficient template processing systems. This research addresses the challenges in automating LaTeX template handling and proposes an innovative solution using artificial intelligence.

Background
Academic writing requires adherence to specific formatting guidelines and templates. Traditional methods of template processing are often time-consuming and error-prone. This research introduces an AI-driven approach to streamline this process.

Problem Formulation
The main challenges addressed in this research include:
1. Efficient template analysis and processing
2. Accurate content generation following academic standards
3. Multilingual support for template handling

Research Objectives
This study aims to:
1. Develop an automated system for LaTeX template processing
2. Optimize template analysis using machine learning
3. Improve the efficiency of academic document preparation

Research Limitations
The scope of this research is limited to:
1. Academic templates in Indonesian and English
2. LaTeX-based document formats
3. Specific academic document types

Benefits of Research
This research provides:
1. Improved efficiency in academic document preparation
2. Reduced errors in template processing
3. Enhanced accessibility for academic writers
4. Time savings for institutions and researchers

\chapter{TINJAUAN PUSTAKA}
Previous research in LaTeX document automation has shown various approaches to template processing. Smith et al. (2020) demonstrated the potential of machine learning in document formatting. Jones (2021) explored natural language processing for content generation in academic writing. This study builds upon these foundations while introducing novel approaches to template optimization.

% Original content (commented out):
% \documentclass{article}
\usepackage[utf8]{inputenc}
\usepackage[margin=0.75in]{geometry}

\begin{document}
	\begin{center}
    
    	% MAKE SURE YOU TAKE OUT THE SQUARE BRACKETS
    
		\LARGE{\textbf{Capstone Proposal}} \\
        \vspace{1em}
        \Large{Complete or Obsolete: A re-evaluation of research} \\
        \vspace{1em}
        \normalsize\textbf{Joseph Eppich} \\
        \normalsize{Eppichj@sunypoly.edu} \\
        \vspace{1em}
        \normalsize{Advisor: Joshua White} \\
        \vspace{1em}
        \normalsize{State University of New York Polytechnic Institute, Utica NY} \\
        \normalsize{Bachelors of Science in Network and Computer Security}
     
	\end{center}
    \begin{normalsize}
    
    	\section{Objective:}
        
        The objective of this capstone is to read the paper "Quantitative Analysis of Intrusion Detection Systems:Snort and Suricata" and attempt at least one test present in the paper to see if the research is still valid today.    \\
      
		\section{Rationale:}
        
       This project is important to me because I believe it is important to always re-evaluate research. With how quickly technology changes research that may have been valid  in 2011 may not longer be valid with today's technology.  The paper i have chosen is about the intrusion detection systems of Suricata and Snort, i hope to get a better understanding of how they work though this capstone and use that knowledge in the near future.
        
	   	\section{Approach:}
        
      I plan on approaching this capstone through several steps.  First I will conduct research on the paper I have chosen, I will do this by reading the paper and taking notes. Second I plan on taking at least one of the tests used in the paper and preforming it myself recording my results.  Lastly I will compare my results with the intials results and record my findings.
        
    	\section{Timeline:}
        
        Give a bulleted timeline of when you will have things completed. Don't forget to include the Student Undergraduate Research Expo poster presentations, and final presentation in your timeline. 
        \\
        
        \noindent You can create bullets by doing:
        
        \begin{itemize}
        
\item(2/25)have a basic idea and outline for my paper
\item (3/1) have my capstone proposal made and submitted, and begin researching
\item(4/19) have my paper finished submitted for final review
\item(5/1) proof read and look over all my material and prepare for my presentation
        \end{itemize}
        
    	\section{Possible Issues:}
        Possible problems I foresee with this capstone is my lack of hardware compared to that used in the original content. beyond that technical issues with suricata and snort could lead to issues in the long run.
\end{normalsize}
  
\end{document}


\end{document}